\section{Appendix: 2023 Exam Bridge}

\begin{exambox}[Read this first --- different instructor, different syllabus]
The 2023 \textit{Information Visualization VO} final exam was set by
\term{Gschwandtner, Aigner \& Miksch} and is a \term{45-question multiple-choice}
paper. Your 2026 course (\term{Manuela Waldner}) uses an open-ended format
(describe / sketch / strategy / critique / lists) plus multiple-choice on
\textit{algorithms}. Treat this appendix accordingly:
\begin{itemize}
  \item \textbf{Part A} re-uses 2023 questions whose topics your main guide
        \emph{already} covers --- good drill on shared fundamentals.
  \item \textbf{Part B} is a compact reference for topics that appear in the
        2023 paper but are \emph{not} emphasised in Waldner's topic list
        (time-oriented data, maps, text vis, the Visual Analytics process,
        legacy multivariate glyphs). These are \textcolor{pitred}{lower
        probability} for your exam --- skim, don't memorise.
\end{itemize}
Conversely, the 2023 paper barely touches what Waldner stresses most
(dimensionality reduction, scagnostics, scalable InfoVis, VIS4ML, sets) --- so
do not treat the old exam as a coverage map.
\end{exambox}

% =====================================================================
\subsection{Part A --- Practice on shared fundamentals}
% =====================================================================
Answers and rationales follow each block. These map to the main-guide sections
shown in brackets.

\begin{qbox}[Graphs, trees \& dynamic layout (\textrightarrow\ \S9)]
\textbf{Q. Force-directed layout.} A force-directed layout algorithm \dots\ (mark all)
\begin{enumerate}[label=\alph*.]
  \item models the graph after a physical system
  \item is a method to reorder rows/columns of a matrix
  \item is a graph-drawing method for node-link diagrams
  \item is used to compute heatmap colours
  \item includes an energy \emph{maximisation} algorithm
\end{enumerate}
\textbf{A: a, c.} It is a spring/charge physical analogy that \emph{minimises}
energy (so e is false --- minimisation, not maximisation).

\medskip
\textbf{Q. Edge design in node-link graphs.} (mark all)
\begin{enumerate}[label=\alph*.]
  \item Polyline edges with multiple bends are allowed, but bends should be minimised.
  \item Lombardi (curved) drawings are preferred by users even with no task advantage.
  \item For directed edges, arrows should not overlap other edges/arrow heads.
  \item Straight-line representations allow bends to reduce crossings.
\end{enumerate}
\textbf{A: a, c.} d is self-contradictory (straight lines have no bends); b
overstates the evidence.

\medskip
\textbf{Q. Treemaps.} A treemap \dots\ (mark all): is non-space-filling / visualises
hierarchies / can be squarified / can be nested / shows vegetation.\\
\textbf{A: visualises hierarchies, squarified, nested.} A treemap \emph{is}
space-filling (containment encodes hierarchy, area encodes the leaf value).

\medskip
\textbf{Q. Dynamic-graph aesthetics.} (mark all): scalability only for
vertices/edges / temporal aliases don't matter / ``reducing cognitive load'' =
less to keep in mind / preserve the mental map.\\
\textbf{A: reduce cognitive load, preserve the mental map.} This is the
\term{mental-map preservation} idea --- see graph morphing / foresighted layout
in \S9.

\medskip
\textbf{Q. Tree encoding.} Which encoding has: parent--child links, ``layered''
\& ``radial'' layouts, no cycles; PRO intuitive/overview, CON large structures
don't fit / aspect ratio?\\
\textbf{A: Node-link (trees).} (Treemap/sunburst/icicle are space-filling;
indentation is the outline form.)
\end{qbox}

\begin{qbox}[Interaction, models \& CMV (\textrightarrow\ \S1, \S6, \S12)]
\textbf{Q. Multi-scale / semantic zoom.} A technique that shows different visual
representations depending on \dots: colour chosen / geometric zoom only / zoom
level \emph{and} available space / overview-vs-detail choice.\\
\textbf{A: ``zoom level and available space'', and the overview/detail option.}
This is \term{semantic zoom} (vs purely geometric).

\medskip
\textbf{Q. Focus+Context usability.} (mark all): overview+detail and fisheye
help specific tasks / zooming hurts performance / F+C needs no prior knowledge /
effectiveness depends on task and goals.\\
\textbf{A: ``help specific tasks'' and ``depends on task and goals''.}

\medskip
\textbf{Q. Munzner's nested model captures \dots} (mark all):
data/task abstraction / domain situation / pattern characteristics / search /
algorithm / visual-encoding \& interaction idiom / image properties.\\
\textbf{A: the four levels --- domain situation, data/task abstraction,
visual-encoding/interaction idiom, algorithm.}

\medskip
\textbf{Q. InfoVis Reference Model --- interaction feeds back into \dots} (mark all):
data transformation / table structure / visual mapping / view transformation.\\
\textbf{A: data transformation, visual mapping, view transformation} --- the
three user-controllable pipeline transforms.

\medskip
\textbf{Q. Coordinated Multiple Views \dots} (mark all): two+ views for one
conceptual entity / combine techniques on different representations / colour for
trends / cannot use brushing\&linking / always semantic zoom / often supported by
brushing\&linking.\\
\textbf{A: ``two+ views for one entity'', ``combine techniques'', ``often
brushing \& linking''.}
\end{qbox}

\begin{qbox}[Perception, encoding \& data (\textrightarrow\ \S2--\S5)]
\textbf{Q. Visual perception is \dots} bottom-up only / combination of bottom-up
and top-down / top-down only. \quad\textbf{A: combination of both.}

\medskip
\textbf{Q. Image of similarity/proximity/closure groupings = ?}
pre-attentive / visual variables / techniques / \term{Gestalt principles}.
\quad\textbf{A: Gestalt principles.}

\medskip
\textbf{Q. Image of colour/orientation/size pop-out = ?}
\term{pre-attentive attributes} / visual variables / Gestalt / techniques.
\quad\textbf{A: pre-attentive attributes.}

\medskip
\textbf{Q. T-shirt sizes are measured on which scale?}
\quad\textbf{A: ordinal} (ordered categories, no fixed interval).

\medskip
\textbf{Q. Expressiveness vs effectiveness.} Given encodings $a$ and $b$:
``$a$ not expressive, $b$ expressive'' / both expressive / ``equally expressive,
$b$ more effective''.\\
\textbf{A (this paper): ``a is not expressive, b is expressive.''} Remember the
two rules: \term{expressiveness} = encode all and only the data's information
(match channel type to attribute type); \term{effectiveness} = most important
attribute \textrightarrow\ most effective channel.

\medskip
\textbf{Q. Anscombe's Quartet illustrates \dots} random variables/different
graphs / goals\&measures / \term{same statistics, different graphs}.
\quad\textbf{A: same statistics, different graphs} --- always plot the data;
summary statistics (mean, variance, correlation) can be identical for very
different distributions.

\medskip
\textbf{Q. Data aggregation / summaries.} (mark all): scatterplots confirm
expected 2D patterns / identical summary stats imply identical data / box-plots
show enough to distinguish raw-data characters / \term{violin plots show more
structure than box plots}.\\
\textbf{A: ``scatterplots confirm patterns'' and ``violin > box for structure''.}
(Box plots can hide multi-modality --- the Anscombe lesson again.)

\medskip
\textbf{Q. Parallel coordinates.} (mark all): better than scatterplots for two
dims / \term{re-ordering axes reveals relations} / arranged in a matrix /
\term{each dimension on one of multiple parallel axes}.\\
\textbf{A: ``re-ordering axes reveals relations'' and ``one axis per dimension''.}
\end{qbox}

% =====================================================================
\subsection{Part B --- Topics in 2023 but not in the main guide}
% =====================================================================
\begin{exambox}[Lower probability]
Everything in Part B reflects the \textit{old} (Aigner/Miksch) syllabus. It is
included for completeness in case a legacy question type is re-used. Prioritise
the main guide first.
\end{exambox}

\subsubsection{Time-oriented data (Aigner et al.)}
A large 2023 theme. Two ways to map the time dimension:
\begin{keybox}[Time \textrightarrow\ space vs time \textrightarrow\ time]
\begin{itemize}
  \item \term{Static (time \textrightarrow\ space)}: time mapped to a spatial
        axis (e.g.\ a line chart, a horizon graph). \textbf{Supports direct
        comparison} between points in time and is \textbf{well suited to
        explorative/analytic tasks}. Good for following trends \emph{and}
        comparing.
  \item \term{Animation (time \textrightarrow\ time)}: time mapped to real
        playback time. \textbf{Simplest} conceptual mapping; \textbf{good for
        following trends and movement}, but \textbf{poor for direct comparison}
        between distant time points and weaker for explorative analysis (the
        viewer must hold previous frames in memory).
\end{itemize}
\end{keybox}
\textbf{Modelling viewpoints} for time: \term{ordered}, \term{branching}
(alternative scenarios / what-ifs), \term{multiple perspectives} (several
interpretations of the same span); time can also be \term{linear vs cyclic} and
\term{time instant vs interval-based}. Mapping time to \emph{visual features}
(e.g.\ colour) is possible but is not direct comparison.
\begin{center}
\begin{tikzpicture}[font=\scriptsize]
  % ---- cyclic: spiral ----
  \begin{scope}
    \draw[->,>=Stealth, thick, tuwblue]
      plot[smooth, domain=0:1080, samples=120, variable=\t]
      ({(0.3+\t/1080*1.0)*cos(\t)},{(0.3+\t/1080*1.0)*sin(\t)});
    \node[below=0.2cm] at (0,-1.4) {cyclic: spiral / ring};
    \node[above] at (0,1.5) {periods align};
  \end{scope}
  % ---- linear axis ----
  \begin{scope}[xshift=5cm]
    \draw[->,>=Stealth, thick, tuwblue] (0,0) -- (3.2,0) node[right] {time};
    \draw plot[smooth] coordinates {(0,0.4)(0.5,0.9)(1,0.3)(1.5,0.9)(2,0.3)(2.5,0.9)(3,0.4)};
    \node[below=0.2cm] at (1.6,-0.3) {linear: straight axis};
  \end{scope}
\end{tikzpicture}
\end{center}
\sketchnote{Cyclic-time techniques wrap the axis into a ring/spiral (e.g.\ a
spiral graph) so periodic patterns line up; linear-time techniques use a
straight axis.}

\subsubsection{Geospatial data \& maps}
\begin{keybox}[Maps cheat-sheet]
\begin{itemize}
  \item \term{Phenomena dimensionality}: point (0-D, e.g.\ a city), line (1-D,
        e.g.\ a street/river), area (2-D, e.g.\ a country region), surface (2.5-D,
        specified by longitude, latitude \emph{and} height).
  \item \term{Choropleth map}: colours predefined regions by a statistical
        value. Assumes the value is (roughly) uniform within each region; better
        than a cartogram at communicating \emph{where} regions are.
  \item \term{Cartogram}: \emph{distorts} region geography so area encodes the
        statistical value --- good for the magnitude, bad for recognisable
        geography.
  \item \term{Colour maps on maps}: sequential for ordered magnitude (e.g.\
        population density); \term{diverging} for a signed change around a
        neutral midpoint (e.g.\ better/worse than last year).
  \item \term{Map projections}: \term{conformal} preserves local angles/shapes;
        \term{equal-area} preserves area; \term{azimuthal} preserves direction
        from a central point; \term{conic} projections make longitude/latitude
        roughly orthogonal. No flat map preserves everything (sphere
        \textrightarrow\ plane is lossy).
\end{itemize}
\end{keybox}

\subsubsection{Text visualization}
\begin{itemize}
  \item \term{Word cloud / tag cloud}: term frequency mapped to font size (does
        \emph{not} convey context/word order).
  \item \term{ThemeRiver / stream graph}: a stacked-area flow showing how
        themes/topics in a document collection change \emph{over time}.
  \item Levels of text representation: \term{lexical} (tokens), \term{syntactic}
        (a token's grammatical function), \term{semantic} (meaning). Text vis is
        useful to compare patterns \emph{across} documents.
\end{itemize}

\subsubsection{Legacy multivariate / glyph techniques}
\begin{keybox}[Glyphs, pixels, SDOF]
\begin{itemize}
  \item \term{Glyph / icon techniques}: built from graphical primitives; encode
        \emph{multiple} data values in one mark's features (size, angle,
        colour\dots). A glyph encodes several values, not one meaning as a whole.
  \item \term{Chernoff faces}: a glyph mapping data dimensions to facial
        features. Suspected to convey emotion; features are \emph{not} perceived
        equally strongly (a known weakness), so they are \emph{not} ideal for
        precise quantitative communication.
  \item \term{Pixel-based techniques}: each data item = at least one pixel,
        attribute value \textrightarrow\ colour map; designed for \emph{large/high}-dimensional
        data, not low-dimensional only.
  \item \term{Semantic Depth of Field (SDOF)}: blurs items by relevance (sharp =
        relevant); can encode an extra dimension in a scatterplot. (A focus+context
        cue, related to \S6.)
\end{itemize}
\end{keybox}

\subsubsection{Visual Analytics process \& design methodology}
\begin{keybox}[VA / methodology one-liners]
\begin{itemize}
  \item \term{Visual Analytics goal}: ``detect the expected and discover the
        unexpected''; synthesise insight from massive, dynamic, ambiguous,
        conflicting data. Components: \term{visualization} + \term{(data)
        mining} + \term{human perception/interaction} in a loop (Keim et al.\
        process: data \textrightarrow\ \{visualization, model\} with feedback).
  \item \term{Knowledge crystallization loop}: first action is \term{forage for
        data} (then structure \textrightarrow\ visualize \textrightarrow\ act).
  \item \term{Visualization design triangle} (Miksch \& Aigner 2014): \term{data},
        \term{users}, \term{tasks}.
  \item \term{Design study methodology} (Sedlmair et al.\ 2012): a vis design
        study fits when the data/task lie in the middle zone --- some required
        information is in the user's head and/or the task is not yet crisply
        defined (not fully automatable, not purely in-the-head).
  \item \term{Guidance}: a computer-assisted process that actively helps resolve
        a \emph{knowledge gap} the user hits during a VA session.
  \item \term{Formative} evaluation = during analysis/design (problem analysis);
        \term{summative} = after implementation.
  \item \term{Problem analysis} asks: who are the users, what data do they work
        with, what are their tasks. \term{User-/human-centred design} benefits:
        higher user productivity, lower training/support cost.
  \item \term{Interaction modelling}: capture requirements to build/evaluate
        against, and validate the effectiveness/guidance of interaction designs.
\end{itemize}
\end{keybox}

\begin{pitfall}[Don't over-invest here]
If time is short, Part B is the first thing to cut. The 2026 exam's own topic
slides do not list time-oriented data, maps, text vis, or the VA process ---
they list dimensionality-reduction algorithms, scalable InfoVis, VIS4ML, sets,
and the strategy/critique/validation question types in \S13.
\end{pitfall}
